%% kapa: exam=1 round=29 year=2018 period=1 form=A subject=real_estate subject_ko=부동산학원론 seq=26 number=106
%% source: Q-Net 감정평가사 기출문제 (cst00309) artlSeq=5202808 file=extracted/1008/5202808/2018년 제29회 감정평가사 1차시험 기출문제/2018년 제29회 감정평가사 1차시험 기출문제/1교시 A형.pdf
%% conversion: kordoc 4.24.51 -> tools/gen1.py (auto); review=proofread-sonnet; flags=none
\begin{question}{106}
부동산관리에 관한 설명으로 옳은 것은?
\begin{choices}
  \item 시설관리(facility management)는 부동산시설의 자산 및 부채를 종합관리하는 것으로 시설사용자나 기업의 요구에 따르는 적극적인 관리에 해당한다.
  \item 자기관리방식은 입주자와의 소통 측면에 있어서 위탁관리방식에 비해 유리한 측면이 있다.
  \item 위탁관리방식은 자기관리방식에 비해 기밀유지가 유리한 측면이 있다.
  \item 혼합관리방식은 자기관리방식에 비해 문제발생시 책임소재 파악이 용이하다.
  \item 건물의 고층화와 대규모화가 진행되면서 위탁관리방식에서 자기관리방식으로 바뀌는 경향이 있다.
\end{choices}
\end{question}
\answer{2}
